\subsection{Module 2: G-Protein Cycle}\label{sec:m2}

\paragraph{Role.}
cAR1 couples to the heterotrimer G$\alpha_2\beta\gamma$.
FRET between the $\alpha$ and $\beta$ subunits shows that the heterotrimer dissociates and reassociates rapidly on addition and removal of cAMP, that during continuous stimulation the activation reaches a dose-dependent steady state and does \emph{not} decline while the physiological responses subside, and that occupied receptors catalyse the G-protein cycle whether or not they are phosphorylated~\cite{janetopoulos2001}.
Adaptation therefore occurs downstream of the G protein, and free G$\beta\gamma$ is the persistent excitatory signal that Module~3 differentiates.
The initial Ras response to uniform cAMP requires G$\beta$ but not G$\alpha_2$~\cite{kortholt2013}, and symmetry breaking and amplification occur between the G protein and Ras~\cite{kataria2013}.
The module follows the G-protein module of Biswas et al.~\cite{biswas2021}.

\begin{table*}[!t]
\centering
\caption{Module~2 reactions: the G-protein cycle of G$\alpha_2\beta\gamma$. G$\alpha_2^{\mathrm{GTP}}$ and G$\alpha_2^{\mathrm{GDP}}$ are the GTP- and GDP-bound free $\alpha$ subunits; G$\beta\gamma$ is the free dimer. A coloured reaction is the one to which the equally coloured statement in the text and the equally coloured passage in the cited paper refer; a small square appends a further statement.}
\label{tab:m2rxn}
\footnotesize
\renewcommand{\arraystretch}{1.15}
\begin{tabularx}{\textwidth}{@{}l>{\raggedright\arraybackslash}p{0.37\textwidth}>{\raggedright\arraybackslash}X@{}}
\toprule
ID & Reaction & Propensity (rate law)\\
\midrule
\mkt{M2-bis-ke0}{2.0} & \mkt{M2-bis-ke0}{$\mathrm{G}_{\alpha\beta\gamma} \to \mathrm{G}\alpha_2^{\mathrm{GTP}}+\mathrm{G}_{\beta\gamma}$} & $k_{E0}[\mathrm{G}_{\alpha\beta\gamma}]$ \\
\mkn{M2-jan-fret}{M2-bis-g}{2.1} & \mkn{M2-jan-fret}{M2-bis-g}{$\mathrm{G}_{\alpha\beta\gamma}+\mathrm{RC} \to \mathrm{G}\alpha_2^{\mathrm{GTP}}+\mathrm{G}_{\beta\gamma}+\mathrm{RC}$} & $k_E[\mathrm{G}_{\alpha\beta\gamma}][\mathrm{RC}]$ \\
\mkn{M2-jan-phos}{M2-bis-eq}{2.2} & \mkn{M2-jan-phos}{M2-bis-eq}{$\mathrm{G}_{\alpha\beta\gamma}+\mathrm{R}_p\mathrm{C} \to \mathrm{G}\alpha_2^{\mathrm{GTP}}+\mathrm{G}_{\beta\gamma}+\mathrm{R}_p\mathrm{C}$} & $\varepsilon_E k_E[\mathrm{G}_{\alpha\beta\gamma}][\mathrm{R}_p\mathrm{C}]$ \\
\mkt{M2-bis-g}{2.3} & \mkt{M2-bis-g}{$\mathrm{G}\alpha_2^{\mathrm{GTP}} \to \mathrm{G}\alpha_2^{\mathrm{GDP}}$} & $k_\mathrm{hyd}[\mathrm{G}\alpha_2^{\mathrm{GTP}}]$ \\
\mkn{M2-tang-reass}{M2-bis-g}{2.5} & \mkn{M2-tang-reass}{M2-bis-g}{$\mathrm{G}\alpha_2^{\mathrm{GDP}}+\mathrm{G}_{\beta\gamma} \to \mathrm{G}_{\alpha\beta\gamma}$} & $k_\mathrm{re}[\mathrm{G}\alpha_2^{\mathrm{GDP}}][\mathrm{G}_{\beta\gamma}]$ \\
\bottomrule
\end{tabularx}
\end{table*}

\begin{table*}[!t]
\centering
\caption{Module~2 constants.}
\label{tab:m2const}
\footnotesize
\renewcommand{\arraystretch}{1.15}
\begin{tabularx}{\textwidth}{@{}l r l c >{\raggedright\arraybackslash}X@{}}
\toprule
Symbol & Value & Unit & Basis & Meaning, justification, source\\
\midrule
$G{:}R$ & \num{60} &  & E & G-protein to receptor ratio; \mkt{M2-bis-ratio}{Biswas et~al. use 3}~\cite{biswas2021}. Chosen to raise the G$\beta\gamma$ count (shot noise) at constant activated fraction; an explicit trade of abundance realism for signal-to-noise, not a claim about abundance. \\
$G_\mathrm{tot}$ & \num{184620} & $\mathrm{molec/voxel}$ & D & $R_\mathrm{tot}\times G{:}R$ ($38\,\mu$M). \\
$t_{1/2,\mathrm{diss}}$ & \num{3.2} & $\mathrm{s}$ & L & Half-time of FRET loss (G-protein dissociation) after stimulation, \mkt{M2-tang-diss}{$3.4\pm1.4$\,s}~\cite{tang2014}; \mkt{M2-bis-tdiss}{reproduced by the model of}~\cite{biswas2021} (3.0--3.3\,s, Table~2). \\
$k_E$ & \num{0.339} & $\mu\mathrm{M}^{-1}\mathrm{s}^{-1}$ & D & $k_E=\ln 2/(t_{1/2,\mathrm{diss}}\,[\mathrm{R}]_\mathrm{tot})$: pseudo-first-order dissociation when all receptors are occupied. \\
$\varepsilon_E$ & \num{1} &  & L & \mkt{M2-jan-phos}{Phosphorylated and unphosphorylated occupied receptors catalyse the G-protein cycle equally}~\cite{janetopoulos2001}; \mkt{M2-bis-eq}{Biswas et~al. likewise weight all occupied classes equally}~\cite{biswas2021}. \\
$k_{E0}$ & \num{2.18e-7} & $\mathrm{s^{-1}}$ & L & Basal, receptor-independent exchange (\mkt{M2-bis-ke0}{Table~2 of}~\cite{biswas2021}, \mkt{M2-bis-ke0}{value for $G{:}R=3$}~\cite{biswas2021}, not rescaled); sets the resting G$\beta\gamma$ floor. \\
$k_\mathrm{hyd}$ & \num{0.162} & $\mathrm{s^{-1}}$ & C & $7\ln2/30$: intrinsic GTP hydrolysis of G$\alpha_2$, raised 7$\times$ above the value matching the 30\,s recovery anchor to de-saturate the stage (see text). \\
$k_\mathrm{re}$ & \num{6.09} & $\mu\mathrm{M}^{-1}\mathrm{s}^{-1}$ & C & $3500/(t_{1/2,\mathrm{re}}\,[\mathrm{G}]_\mathrm{tot}/2)$ with $t_{1/2,\mathrm{re}}=30$\,s (measured reassociation half-times \mkt{M2-tang-reass}{28.7--32.3\,s}~\cite{tang2014}); the factor 3500 keeps recapture pseudo-first order (see text). The realised reassociation half-time is 7.5\,s, not 30\,s. \\
$D_{\mathrm{G}\alpha\beta\gamma},D_{\mathrm{G}\alpha_2},D_{\mathrm{G}\beta\gamma}$ & \num{0.2} & $\mu\mathrm{m^2/s}$ & E & The value $0.2\,\mu\mathrm{m}^2/\mathrm{s}$ is that of \mkt{M2-bis-D}{Table~3 of}~\cite{biswas2021}; equal mobility of the three forms is a design choice (G proteins also \mkt{M2-elzie}{exchange between cytosol and membrane}~\cite{elzie2009}). \\
\bottomrule
\end{tabularx}
\end{table*}


\paragraph{Reactions.}
\begin{description}\itemsep0pt
\item[2.0] Basal, receptor-independent nucleotide exchange. Its flux is negligible (about 0.03 molecules\,s$^{-1}$ per cell) but it defines the resting G$\beta\gamma$ floor, so that the receiver has a defined baseline rather than a hard zero.
\item[2.1] Receptor-catalysed exchange: an occupied receptor dissociates the heterotrimer into G$\alpha_2$-GTP and G$\beta\gamma$ and is not consumed.
\item[2.2] The same reaction for the phosphorylated occupied receptor. Both forms catalyse the cycle with equal efficiency ($\varepsilon_E=1$), as observed~\cite{janetopoulos2001}, so the phosphorylated receptor carries the drive at high cAMP.
\item[2.3] Intrinsic GTP hydrolysis converts G$\alpha_2$-GTP into G$\alpha_2$-GDP. An RGS-accelerated hydrolysis (2.4) is defined in the code but switched off.
\item[2.5] Recapture of G$\beta\gamma$ by G$\alpha_2$-GDP re-forms the heterotrimer and terminates the signal.
\end{description}

\paragraph{Constants.}
$k_E$ is derived from the dissociation half-time at saturating cAMP: when all receptors are occupied the heterotrimer is consumed with the pseudo-first-order rate $k_E[\mathrm R]_\mathrm{tot}$, so $k_E=\ln2/(t_{1/2,\mathrm{diss}}[\mathrm R]_\mathrm{tot})$ with $t_{1/2,\mathrm{diss}}=3.2$\,s, the half-time of FRET loss measured in living cells, $3.4\pm1.4$\,s~\cite{tang2014}, which the model of Biswas et al.\ reproduces~\cite{biswas2021}.
The basal constant $k_{E0}$ is the value of Table~2 of~\cite{biswas2021} for a ratio of 3.

\emph{Why the gain of the stage was changed.}
The conservation law $[\mathrm G\beta\gamma]=[\mathrm G\alpha_2^\mathrm{GTP}]+[\mathrm G\alpha_2^\mathrm{GDP}]$ holds exactly, because every dissociation creates one G$\alpha_2$ and one G$\beta\gamma$ and every recapture removes one of each.
In the steady state, with $P$ the total exchange flux, $[\mathrm G\alpha_2^\mathrm{GTP}]=P/k_\mathrm{hyd}$ and $k_\mathrm{re}[\mathrm G\alpha_2^\mathrm{GDP}][\mathrm G\beta\gamma]=P$.
If recapture is slow compared with hydrolysis, $[\mathrm G\beta\gamma]\approx\sqrt{P/k_\mathrm{re}}$ and the stage compresses the input contrast (exponent $\tfrac12$); if recapture is fast, $[\mathrm G\beta\gamma]\approx[\mathrm G\alpha_2^\mathrm{GTP}]=P/k_\mathrm{hyd}$ is linear in the drive until the heterotrimer is depleted.
The exponent is therefore a property of the regime and not of the conservation law.
Raising $k_\mathrm{re}$ by a factor of 3500 puts the stage in the linear regime, and raising $k_\mathrm{hyd}$ by a factor of 7 moves the saturation point to higher cAMP.
On the stochastic engine, with replicate cells, the contrast transfer $\mathrm d\ln\mathrm G\beta\gamma/\mathrm d\ln c$ at 60\,nM cAMP increased from 0.22 to 0.77 and the signal-to-noise ratio doubled; at 5\,nM the transfer increased from 0.85 to 1.65 with no loss of signal-to-noise ratio.

\emph{Departures from the literature anchors.}
The chosen values move the two half-time anchors: the dissociation half-time becomes about 2\,s (from 3.2\,s) and the reassociation half-time about 7.5\,s, four times shorter than the measured 29--32\,s~\cite{tang2014}.
The cost of the stochastic simulation at the working point rises about threefold, because the activated fraction is increased on purpose.
The G-protein pool is 60 times the receptor number, 20 times the ratio 3 of~\cite{biswas2021}.
This pool size raises the G$\beta\gamma$ copy number at a fixed activated fraction and was accompanied by paired changes in Module~3 (see there) that leave its dynamics unchanged; it is a deliberate trade of realism for signal-to-noise and is not a claim about G-protein abundance.
The three G-protein forms have the same diffusion constant, which is a design choice.
